\documentclass[10pt]{article}

% ---------- Document Identity (update these only) ----------
\newcommand{\DocStage}{}
\newcommand{\DocTitle}{Your Road to Tier One SOF}
\newcommand{\ProgramName}{182-Week Civilian-to-Selection Readiness Pipeline}
\newcommand{\ProgramSubtitle}{}

% ---------- Page ----------
\usepackage[legalpaper,left=0.5in,right=0.5in,top=0.6in,bottom=0.45in]{geometry}

% ---------- Typography ----------
\usepackage[T1]{fontenc}
\usepackage[utf8]{inputenc}
\usepackage{libertinus}
\usepackage{microtype}
\usepackage{parskip}
\usepackage{ragged2e}
\usepackage{textcomp}
\AtBeginDocument{\color{black}}
\AtBeginDocument{\pagecolor{white}}
\AtBeginDocument{\justifying}

% ---------- Landscape pages ----------
\usepackage{pdflscape}

% ---------- Color ----------
\usepackage[table]{xcolor}
\definecolor{StageBlue}{HTML}{1F4E79}
\definecolor{StageBlueDark}{HTML}{163A5A}
\definecolor{LightGray}{HTML}{F2F4F7}
\definecolor{MidGray}{HTML}{D9DEE5}
\definecolor{RuleGray}{HTML}{C7CED8}
\definecolor{HeaderInk}{HTML}{000000}
\definecolor{Ink}{HTML}{000000}

% ---------- Utility ----------
\usepackage{enumitem}
\sloppy
\setlength{\emergencystretch}{2em}

% ---------- Boxes / UI ----------
\usepackage[most]{tcolorbox}

\tcbset{
  charterTitle/.style={
    enhanced,
    colback=StageBlue!4,
    colframe=StageBlue,
    boxrule=1.25pt,
    arc=3pt,
    left=16pt,right=16pt,top=14pt,bottom=14pt,
    borderline north={3.2pt}{0pt}{StageBlue},
    borderline west={4pt}{0pt}{StageBlue},
    borderline south={1.25pt}{0pt}{StageBlue},
    drop shadow={black!25!white},
  },
  charterRule/.style={
    enhanced,
    colback=LightGray,
    colframe=RuleGray,
    boxrule=0.85pt,
    arc=3pt,
    left=12pt,right=12pt,top=10pt,bottom=10pt,
    borderline west={3pt}{0pt}{StageBlue},
  },
  charterBadge/.style={
    enhanced,
    colback=StageBlue,
    coltext=white,
    colframe=StageBlueDark,
    boxrule=0.6pt,
    arc=2pt,
    left=6pt,right=6pt,top=3pt,bottom=3pt,
  }
}
\newtcbox{\Badge}{nobeforeafter,charterBadge}

% ---------- Lists ----------
\setlist[itemize]{leftmargin=1.5em, itemsep=0.25em, topsep=0.25em}
\setlist[enumerate]{leftmargin=1.8em, itemsep=0.25em, topsep=0.25em}

% ---------- TOC ----------
\usepackage{tocloft}
\setcounter{tocdepth}{3}
\setlength{\cftbeforesecskip}{3pt}
\setlength{\cftbeforesubsecskip}{2pt}
\setlength{\cftbeforesubsubsecskip}{0pt}
\renewcommand{\cftsecfont}{\bfseries}
\renewcommand{\cftsecpagefont}{\bfseries}
\renewcommand{\cftsubsecfont}{\normalfont}
\renewcommand{\cftsubsecpagefont}{\normalfont}
\renewcommand{\cftsubsubsecfont}{\normalfont}
\renewcommand{\cftsubsubsecpagefont}{\normalfont}
\setlength{\cftsecindent}{0pt}
\setlength{\cftsubsecindent}{1.25em}
\setlength{\cftsubsubsecindent}{2.8em}
\setlength{\cftsecnumwidth}{2.1em}
\setlength{\cftsubsecnumwidth}{2.8em}
\setlength{\cftsubsubsecnumwidth}{3.6em}

% Remove the built-in TOC heading so only the custom heading is shown
\makeatletter
\renewcommand{\tableofcontents}{\@starttoc{toc}}
\makeatother

% ---------- Header/Footer: page number only ----------
\usepackage{fancyhdr}
\pagestyle{fancy}
\fancyhf{}
\renewcommand{\headrulewidth}{0pt}
\renewcommand{\footrulewidth}{0pt}
\fancyfoot[C]{\thepage}

% ---------- Section formatting ----------
\usepackage{titlesec}

\titleformat{\section}
  {\Large\bfseries\color{Ink}}
  {\thesection}{0.6em}{}
  [\vspace{0.25em}\textcolor{StageBlue}{\rule{\linewidth}{0.8pt}}]

\titleformat{\subsection}
  {\large\bfseries\color{Ink}}
  {\thesubsection}{0.6em}{}

\titleformat{\subsubsection}
  {\normalsize\bfseries\color{Ink}}
  {\thesubsubsection}{0.6em}{}

\titlespacing*{\section}{0pt}{0.95em}{0.35em}
\titlespacing*{\subsection}{0pt}{0.65em}{0.25em}
\titlespacing*{\subsubsection}{0pt}{0.55em}{0.2em}

% ---------- Table engine ----------
\usepackage{tabularray}
\UseTblrLibrary{booktabs}

% ---------- tabularray: GLOBAL table treatment ----------
\SetTblrInner{
  cells = {fg=black, bg=white, valign=t},
}

% ---------- Header helpers ----------
\newcommand{\Hdr}[1]{\shortstack[c]{\strut #1\strut}}

% ---------- Lines ----------
\newcommand{\TitleLine}{\textcolor{StageBlue}{\rule{0.98\linewidth}{0.95pt}}}
\newcommand{\SubTitleLine}{\textcolor{RuleGray}{\rule{0.98\linewidth}{0.65pt}}}

% ---------- Continued on next page ----------
\newcommand{\ContinuedOnNextPageInline}{%
  \par
  \vfill
  \noindent\textcolor{RuleGray}{\rule{\linewidth}{1pt}}\vspace{-2pt}
  \noindent\hfill\textit{Continued on next page}\par\vspace{-10pt}
  \noindent\textcolor{RuleGray}{\rule{\linewidth}{1pt}}
  \clearpage
}

% ---------- PDF hyperlinks ----------
\usepackage[hidelinks]{hyperref}
\hypersetup{
  colorlinks=false,
  pdfborder={0 0 0},
  linktoc=all,
  pdftitle={\DocTitle},
  pdfauthor={\ProgramName}
}

% =========================================================
\begin{document}

\section*{7.09 — Health Monitoring, Diagnostic-Adjacent Monitoring, and Biometrics}
\phantomsection
\addcontentsline{toc}{section}{7.09 — Health Monitoring, Diagnostic-Adjacent Monitoring, and Biometrics}

Overview \textbf{ID}: 7.09\\
Overview Name: Health Monitoring, Diagnostic-Adjacent Monitoring, and Biometrics\\
Version: v1.0\\
Governing Status: Draft — Non-Deployable\\
Scope Boundary Statement: Covers non-diagnostic monitoring tools and routines for trend visibility, conservative safety-signal capture, and evidence continuity. Does not cover diagnosis, treatment, lab ordering, clinical interpretation, replacement of clinician judgment, or Testing, Timing, Measurement, and Course-Setup Tools (7.18).

\subsection*{7.09.1 Purpose and Scope}
\phantomsection
\addcontentsline{toc}{subsection}{7.09.1 Purpose and Scope}

Purpose:

7.09 defines non-diagnostic monitoring tools and routines that support readiness-trend and recovery-trend visibility, safety-signal detection through conservative, standardized measurement capture, and audit-grade evidence continuity and stability through standardized routines and tool-identity logging.

In scope:

\begin{itemize}
\item baseline biometric and illness-context tools (scale, tape, and basic thermometer),
\item optional consumer wearable trackers,
\item optional chest-strap heart-rate monitors, and
\item optional diagnostic-adjacent tools only in clinician-guided, record-only posture.
\end{itemize}

Out of scope:

\begin{itemize}
\item medical diagnosis or treatment,
\item lab ordering,
\item clinical interpretation,
\item replacement of clinician judgment, and
\item Testing, Timing, Measurement, and Course-Setup Tools (7.18).
\end{itemize}

This overview does not prescribe workouts, progression rules, testing protocols, or treatment plans.

No device readout, recovery score, or trend output in this overview constitutes clearance to continue training when symptoms or stop rules are present.

Diagnostic-adjacent values may be logged only; they may not be self-interpreted or used to drive self-directed treatment decisions.

\subsection*{7.09.2 Governance and Safety Boundaries}
\phantomsection
\addcontentsline{toc}{subsection}{7.09.2 Governance and Safety Boundaries}

\begin{itemize}
\item Monitoring is decision support, not diagnosis, not treatment advice, not a substitute for clinician evaluation, and not clearance to continue training when symptoms are present.
\item The decision hierarchy for this category is controlling:
  \begin{itemize}
  \item symptoms, stop rules, and acute safety concerns override device readouts, trend outputs, and summary scores;
  \item materially concerning symptoms plus materially concerning readings trigger immediate stop and clinician evaluation (non-diagnostic).
  \end{itemize}
\item Trend summaries, recovery scores, and device alerts must never override symptoms, stop rules, or acute safety concerns.
\item When materially concerning symptoms or readings are present, stop training and route the case to clinician evaluation rather than continuing data collection.
\item Repeated self-checks must not be used to delay clinician evaluation when stop rules or materially concerning symptoms are present.
\item Consistency requirements:
  \begin{itemize}
  \item the same device, the same measurement conditions, the same time of day where relevant, and the same pre-measurement state;
  \item the same fasted/fed, rested, and hydration-state conditions where relevant;
  \item consistent rounding precision and unit conventions across all logged entries;
  \item consistent anatomical landmarks and tape tension for tape-measure sites.
  \end{itemize}
\item When a monitoring tool changes, log the tool change and treat the trend as re-baselined until comparability is re-established.
\item Missing entries must be recorded explicitly as missing or unknown and must not be reconstructed from memory.
\item Do not mix units or rounding conventions within an active trend line unless the change is logged and the series is re-baselined.
\item If anatomical landmarking changes for tape-based tracking, log the change and restart the trend rather than blending unlike measurements.
\item No optimization claims: this category is for visibility, trend continuity, and audit stability, not self-optimization or performance-manipulation practices.
\item No device output in this category may be used to justify overriding stop rules, acute symptoms, or conservative reslot decisions.
\item Consumer-device readiness features may be logged for trend visibility only and must not be framed as optimization prescriptions.
\item Diagnostic-adjacent tools are optional, clinician-guided, and record-only:
  \begin{itemize}
  \item record-only; no self-interpretation or self-directed action;
  \item do not use them to self-manage symptoms, diagnoses, treatment decisions, or training clearance.
  \end{itemize}
\item Materially concerning values must be logged and routed to clinician evaluation rather than interpreted inside this overview.
\item When clinician-guided diagnostic-adjacent monitoring is active, log the tool, the date introduced, and the record-only posture in the evidence file.
\end{itemize}

\subsection*{7.09.3 Tier Doctrine — Applied to Monitoring and Biometrics}
\phantomsection
\addcontentsline{toc}{subsection}{7.09.3 Tier Doctrine — Applied to Monitoring and Biometrics}

\begin{itemize}
\item Tier 0: self-report daily-readiness fields, symptom-screen observation, and minimal manual measures (including resting manual pulse count), strictly non-diagnostic and record-only.
\item Missing Tier 0 entries must be marked explicitly as missing or unknown and must not be estimated from memory.
\item Tier 1: minimum deployable measurement-tool set for baseline biometric trends and conservative illness-context capture (scale, tape, and basic thermometer), with standardized routines, stable measurement conditions, and logged tool identity.
\item Tier 2: expanded tools for better trend visibility and session-intensity logging (wearable HR/sleep tracker, chest-strap heart-rate monitor, pulse oximeter in clinician-guided record-only posture, and upper-arm blood-pressure cuff in record-only posture), still non-diagnostic.
\item Additional devices increase trend visibility only; they do not create diagnostic authority, treatment guidance, or self-clearance authority.
\item Tier 2 devices must be logged under consistent wear, placement, and timing conditions to preserve comparability.
\item Tier 3: redundancy and higher-fidelity trend capture (heart-rate-variability-capable wearable, backup scale and backup tape, backup upper-arm blood-pressure cuff where clinician-guided blood-pressure capture is active, and offline log-backup posture), emphasizing stability, reduced variance, and evidence continuity rather than expanded metric count.
\item Tool switches must be logged and treated as re-baselining events until comparability is re-established.
\item Offline backup is an evidence-continuity control and does not constitute a new measurement category.
\item Backup tools must be checked for reasonable comparability before they replace primary logging tools in an active trend line.
\end{itemize}

\subsection*{7.09.4 Selection Criteria \& Fit Checks}
\phantomsection
\addcontentsline{toc}{subsection}{7.09.4 Selection Criteria \& Fit Checks}

Measurement quality criteria:

\begin{itemize}
\item Repeatability and stable routines are prioritized over expanded sensor count.
\item Prefer tools with clear unit display and stable output under consistent conditions.
\item Prefer data-export capability to preserve audit traceability.
\end{itemize}

Fit checks (operational):

\begin{itemize}
\item Blood-pressure cuff sizing (upper-arm cuff) must match arm size; incorrect sizing undermines repeatability.
\item Wearable fit: comfortable for overnight wear; no skin irritation; stable sensor contact.
\item Chest-strap fit: stable contact without chafing; tolerable during longer sessions.
\end{itemize}

Data integrity:

\begin{itemize}
\item Timestamps are captured consistently.
\item Units are consistent and not silently changed.
\item Tool changes are treated as comparability events and logged explicitly.
\end{itemize}

\subsection*{7.09.5 Setup, Safety, and Anchoring}
\phantomsection
\addcontentsline{toc}{subsection}{7.09.5 Setup, Safety, and Anchoring}

Standardization routines:

\begin{itemize}
\item Time-of-day consistency for resting measures where relevant.
\item Rested versus post-exercise distinction is explicit in logs (do not mix contexts).
\item Repeat readings when a reading seems erratic; do not interpret; record the repeat status.
\end{itemize}

Safety posture:

\begin{itemize}
\item Symptoms override device readouts. If a red-flag symptom exists, stop immediately and seek clinician evaluation regardless of normal numbers.
\item Do not use good numbers to justify testing through illness, pain drift, or unusual breathlessness.
\end{itemize}

Privacy and data handling:

\begin{itemize}
\item Record the minimum necessary fields to support readiness and audit.
\item Store logs in the program’s system-of-record posture; avoid fragmented, unsynced sources.
\end{itemize}

\subsection*{7.09.6 Maintenance, Replacement, and Storage}
\phantomsection
\addcontentsline{toc}{subsection}{7.09.6 Maintenance, Replacement, and Storage}

\begin{itemize}
\item Batteries and charging:
  \begin{itemize}
  \item replace batteries as needed and maintain charging routines to prevent missing-data gaps.
  \end{itemize}
\item Cleaning:
  \begin{itemize}
  \item wearable bands or straps and chest straps are cleaned to reduce skin irritation and device failure.
  \item cuffs are cleaned per basic hygiene posture.
  \end{itemize}
\item Non-clinical accuracy posture:
  \begin{itemize}
  \item if outputs become erratic under stable conditions, treat as tool drift and replace or remove from use.
  \end{itemize}
\item Storage:
  \begin{itemize}
  \item dry storage protected from heat and humidity.
  \item prevent damage to tape integrity and cuff tubing.
  \end{itemize}
\end{itemize}

\subsection*{7.09.7 Substitution Rules — Maintain Monitoring Intent}
\phantomsection
\addcontentsline{toc}{subsection}{7.09.7 Substitution Rules — Maintain Monitoring Intent}

\begin{itemize}
\item If a wearable is unavailable:
  \begin{itemize}
  \item use rating of perceived exertion (RPE), talk test, and self-report logs; program validity is preserved because these are governance-level anchors.
  \end{itemize}
\item If a diagnostic-adjacent tool is absent:
  \begin{itemize}
  \item rely on symptom screening and stop rules; do not infer missing metrics.
  \end{itemize}
\item Tool change discipline (Stage 2.E posture, high level):
  \begin{itemize}
  \item record tool-change event (old tool $\rightarrow$ new tool),
  \item treat trends as re-baselined after tool change,
  \item record missing data as missing/unknown; no inference or memory backfill.
  \end{itemize}
\item If deviations affect admissibility:
  \begin{itemize}
  \item apply the appropriate validity tag posture (session/test admissibility remains governed elsewhere); do not claim gate credit on non-comparable data.
  \end{itemize}
\end{itemize}

\subsection*{7.09.8 Phase Integration Map — Required}
\phantomsection
\addcontentsline{toc}{subsection}{7.09.8 Phase Integration Map — Required}

\begingroup
\scriptsize
\setlength{\tabcolsep}{2pt}
\renewcommand{\arraystretch}{1.12}

\begin{longtblr}{
  width=\linewidth,
  colspec = {
    Q[l,wd=1.05in]
    Q[l,wd=1.55in]
    X[l,2.75]
    X[l,3.65]
  },
  rowhead=1,
  hlines,
  vlines,
  cells = {hsep=2pt, vsep=6pt, halign=l, valign=t},
  row{odd} = {bg=white},
  row{even} = {bg=LightGray},
  row{1} = {bg=StageBlue!15, fg=black, font=\bfseries, halign=l, valign=m},
}
\Hdr{Phase} &
\Hdr{Required Tier (from Stage 6)} &
\Hdr{What Changes / Becomes Mandatory} &
\Hdr{Notes} \\
Phase 00 &
T0 &
None &
Optional only. Self-report and symptom screening are still recommended as conservative safety practice but are not tier-required. \\
Phase 01 &
T0 &
None &
Optional only. Measurement discipline should not drift; if tools are used, stabilize routines outside protected windows. \\
Phase 02 &
T0 &
None &
Optional only. If any phase artifact treats 7.09 as required in Phase 02, Requires Stage 8 review. \\
Phase 03 &
T0 &
None &
Optional only. Increasing consequence raises the value of standardization, but Stage 6 does not require Tier 1 yet. \\
Phase 04 &
T0 &
None &
Optional only. Aquatic relevance does not change 7.09 requirements. \\
Phase 05 &
T0 &
None &
Optional only. Endurance consequence increases the value of stable routines; still optional by Stage 6. \\
Phase 06 &
T1 (Entry) &
Tier 1 becomes minimum deployable by Phase 06 (Entry): scale + tape + thermometer routines standardized &
Phase 06 (Entry) is first required point. Establish stable time-of-day routines before protected windows; treat tool changes as re-baselining events. \\
Phase 07 &
T1 (Entry) &
Maintain Tier 1 routines; Tier 2 tools remain optional &
Hybrid density increases variance costs; do not add novelty near protected windows. \\
Phase 08 &
T1 (Entry) &
Maintain Tier 1 routines; emphasize durability trend awareness &
Use logs to enforce conservative downshift/reslot posture; no clinical interpretation. \\
Phase 09 &
T1 (Entry) &
Maintain Tier 1 routines; prevent missing-data drift &
Missing data defaults to missing/unknown; do not infer. \\
Phase 10 &
T1 (Entry) &
Maintain Tier 1 routines; optional Tier 2 may improve trend visibility &
If adding a wearable tracker or chest-strap heart-rate monitor, stabilize outside protected windows; do not treat as required. \\
Phase 11 &
T1 (Entry) &
Maintain Tier 1 routines; optional Tier 2 may support recovery sensitivity &
Symptoms override readouts; do not test through illness. \\
Phase 12 &
T1 (Entry) &
Maintain Tier 1 routines; Tier 3 redundancy is optional only &
Stage 6.E has flagged internal overbuild-then-reduce concerns in Stage 6.B; Requires Stage 8 review if Tier expectations drift. \\
Phase 13 &
T1 (Entry) &
Maintain Tier 1 routines; consolidation emphasizes low variance &
If any artifact implies Tier 3 is required here due to late-stage verification posture, Requires Stage 8 review. \\
\end{longtblr}

\endgroup

7.09.8 Phase Integration Map Notes:

\begin{itemize}
\item Stage 6 timing used here: 7.09 first becomes minimum deployable Tier 1 at Phase 06 (Entry). If any phase demand posture requires 7.09 earlier, Requires Stage 8 review.
\item If Stage 6 internal tier consistency creates a downstream deployability conflict, Requires Stage 8 review.
\end{itemize}

\subsection*{7.09.9 Risks, Contraindications, and Red Flags}
\phantomsection
\addcontentsline{toc}{subsection}{7.09.9 Risks, Contraindications, and Red Flags}

\begin{itemize}
\item False reassurance risk: normal numbers do not override red-flag symptoms; stop posture remains controlling.
\item Anxiety from numbers: over-monitoring can increase anxiety and drive poor decisions; prioritize stable routines and minimal necessary measures.
\item Device error and drift: consumer tools drift; treat erratic readings as tool drift, not as medical truth; replace or omit.
\item Overreacting to single readings: trend visibility is the goal; persistent abnormal readings with symptoms warrant clinician evaluation.
\item Clinician referral triggers (non-diagnostic):
  \begin{itemize}
  \item persistent concerning symptoms,
  \item repeated abnormal readings combined with symptoms,
  \item suspected sleep-apnea signs,
  \item recurrent panic or anxiety spikes,
  \item repeated near-syncope, chest pressure, disproportionate dyspnea.
  \end{itemize}
\item Red-flag symptoms reminder:
  \begin{itemize}
  \item chest pain or pressure, syncope or near-syncope, severe or disproportionate shortness of breath, new neurologic symptoms $\rightarrow$ stop immediately and seek clinician evaluation.
  \end{itemize}
\item Requires Stage 8 review triggers:
  \begin{itemize}
  \item any artifact requiring 7.09 earlier than Phase 06 (Entry),
  \item any artifact using 7.09 tools to justify ignoring symptoms or stop posture,
  \item any artifact adding environmental monitors into 7.09 scope.
  \end{itemize}
\end{itemize}

\end{document}